\documentclass[11pt]{article}
\usepackage[T1]{fontenc}
\usepackage[utf8]{inputenc}
\usepackage{lmodern,amsmath,amssymb,amsthm}
\usepackage[margin=1in]{geometry}
\newtheorem{theorem}{Theorem}
\newcommand{\Z}{\mathbb Z}\newcommand{\C}{\mathbb C}
\newcommand{\E}{\mathbb E}
\title{The plain additive-quadruple implication on a finite interval}
\author{Author proposal: /root/astra\_v14\_baseline\_theta}
\date{7 November 2026}
\begin{document}\maketitle

\noindent\textbf{Scope.} This local proof supplies the exact plain
correlation statement of GTZ Proposition 8.4, p.~1277, needed by the
Section 13 join. It does not use Proposition 8.3 or infer plain bias
from lower-degree bias. The argument is the elementary two-step
Cauchy--Schwarz pattern; the finite domains and constants are written
explicitly. No source error in Proposition 8.4 is asserted. This
author packet requires independent review.

\section{Precise finite assertion}
Put $J=\{1,\ldots,N\}$, $I=\{-N,\ldots,N\}$, and suppose
$H\subseteq I$ has $|H|\geq\alpha |I|$ for fixed $\alpha>0$.
Let $f:J\to\C$ be $1$-bounded and extend it by zero to $\Z$.
For each $h\in H$ let $\chi_h:\Z\to\C$ be globally bounded
by a fixed $B>0$, and assume
\[
 \left|\E_{n\in J} f(n+h)\overline{f(n)}\chi_h(n)\right|
 \geq\delta>0.
\]
The values outside $J$ are part of the data and are not truncated.

\begin{theorem}\label{thm:plain}
Write $a=\alpha\delta/3$ and $\eta=a^4/4$.
There are at least
\[
 \frac{\eta}{2B^4}(4N+1)(2N+1)^2
\]
distinct additive quadruples $(h_1,h_2,h_3,h_4)\in H^4$,
$h_1+h_2=h_3+h_4$, for which
\[
 \left|\E_{n\in J}
  \chi_{h_1}(n)\chi_{h_2}(n+h_1-h_4)
  \overline{\chi_{h_3}(n)}
  \overline{\chi_{h_4}(n+h_1-h_4)}\right|
 \geq\eta/2.
\]
In particular their number is a positive fixed multiple of $N^3$.
\end{theorem}

\section{Two exact Cauchy--Schwarz steps}
For each $h\in H$ choose a scalar $c_h$ of modulus one so the
given average multiplied by $c_h$ is real and at least $\delta$.
Set $u_h=c_h\chi_h$ on $H$, and $u_h=0$ for every $h\notin H$.
The functions remain defined on all integers. Then
\[
 \E_{h\in I}\E_{n\in J}
 f(n+h)\overline{f(n)}u_h(n)\geq\alpha\delta.
\]
Let $M=J+I=\{1-N,\ldots,2N\}$, so $|M|=3N$.
Set $m=n+h$. The preceding average is exactly
\[
 3\E_{m\in M}f(m)
       \E_{h\in I}\overline{f(m-h)}u_h(m-h).
\]
Terms with $m-h\notin J$ vanish because $f$ was zero-extended.
Cauchy--Schwarz in $m$, using $|f|\leq1$, therefore yields
\begin{equation}\label{eq:first}
 \E_{m\in M}\left|\E_{h\in I}
          \overline{f(m-h)}u_h(m-h)\right|^2\geq a^2.
\end{equation}

Expand the square and put $h'=h+k$, $n=m-h$.
Let $K=\{-2N,\ldots,2N\}$. The expanded left side equals
\[
 \frac{N}{|M|}\frac{|K|}{|I|}
 \E_{k\in K}\E_{n\in J}
 \overline{f(n)}f(n-k)T(k,n),
 \qquad
 T(k,n)=\E_{h\in I}u_h(n)\overline{u_{h+k}(n-k)}.
\]
This is an equality with the following domain checks. The factor
$\overline{f(n)}$ forces $n\in J$, and then $m=n+h\in M$
automatically for $h\in I$. The condition $h+k\in I$ is supplied
by $u_{h+k}=0$ outside $H\subseteq I$. Thus no additional cutoff
in $n-k$ is inserted into $T$. Its values there are legitimate because
the original functions are globally defined. The remaining factor
$f(n-k)$ stays in the bounded weight.

The prefactor is positive and at most two. By \eqref{eq:first},
the absolute value of the displayed average without prefactor is
at least $a^2/2$. Cauchy--Schwarz in $(k,n)$, dropping only the
bounded weight $\overline{f(n)}f(n-k)$ by an inequality, gives
\[
 \E_{k\in K}\E_{n\in J}|T(k,n)|^2\geq a^4/4=\eta.
\]
Expand this square with independent $h,h'\in I$. It is the average
over $(k,h,h')\in K\times I^2$ of
\[
 C(k,h,h')=\E_{n\in J}
 u_h(n)u_{h'+k}(n-k)
       \overline{u_{h'}(n)}\overline{u_{h+k}(n-k)}.
\]
The average is real and at least $\eta$, while $|C|\leq B^4$.
Consequently at least a fraction $\eta/(2B^4)$ of triples have
$|C|\geq\eta/2$. Indeed the complementary contribution to the
average of $|C|$ is at most $\eta/2$, and the contribution of the
selected fraction is at most that fraction times $B^4$.

Such a triple has $h,h',h+k,h'+k\in H$, since otherwise $C=0$.
Set
\[
 (h_1,h_2,h_3,h_4)=(h,h'+k,h',h+k).
\]
It is an additive quadruple and $h_1-h_4=-k$. The map is injective,
with inverse $h=h_1$, $h'=h_3$, $k=h_4-h_1$.
The factors $c_h$ contribute only a constant phase of modulus one
to $C$, so they do not change its absolute value. This proves the
theorem, including the exact unweighted average on $J$.

\section{Vector families, internal limits and actual consumer}
For vector-valued $\chi_h$ of fixed finite dimension $D$, a positive
Euclidean correlation implies that some component has magnitude at
least $\delta/\sqrt D$. Finite pigeonhole retains a component on
a subfamily of density at least $\alpha/D$. Apply the scalar theorem
there. The scalar output is a component of the full tensor product,
so it implies plain vector bias. A finite list of dimensions can first
be padded into one fixed dimension. For a bounded $f$ with a fixed
bound other than one, divide it by that bound, changing $\delta$
by a fixed factor only.

Apply the finite assertion to representatives of an internal family,
with uniform standard bounds, density and correlation. Choices of
components and phases can be made for each finite family and then
passed to the ultralimit. All constants above are positive standard
numbers, so the selected quadruples have positive standard density
on the $N^3$ scale and the correlation is positive by a standard
amount. This is exactly the plain Proposition 8.4 output.

In Section 13 the input to this theorem is the entire bounded product
that includes the selected lower nilcharacter. Keeping it inside
$\chi_h$ throughout the calculation preserves that lower information.
The theorem requires no degree bound and introduces no lower-degree
test factor. It therefore supplies the stronger plain statement that
cannot be replaced merely by the degree-biased LP03 output.

The bounded application audit checks the original $J$, the enlarged
$M=J+I$, the parameter range $K$, the exact inverse reindexing and
the global domain of $\chi_h$. A possible artificial cutoff
$1_J(n-k)$ is unnecessary: $f(n-k)$ is removed by the second
Cauchy--Schwarz inequality, not by claiming an equality after erasing
its support. This resolves the domain question without classifying
the published Proposition 8.4 as false or defective. The modified
Proposition 7.3, its nilsequence-specific factorization and the
remaining Section 13 hypotheses are outside this theorem's scope.

\end{document}
